% TOP_PROBLEM: {"id": 1920, "title": "Erdős Problem #86", "category_id": 3, "category": {"id": 3, "name": "graph_theory", "display_name": "Graph Theory", "description": "Problems involving graphs, networks, and their properties.", "slug": "graph-theory", "order_index": 3, "created_at": "2026-07-31T15:26:25.671Z"}, "status": "open"}
% FIRST_POSTED: null
% ENTRY_KIND: statement_only
% Imported from: unsolvedmath_additional_500.tex; UnsolvedMath ID 1920
% !TeX program = lualatex
% Compile twice: lualatex 1920.tex
% Self-contained: no external bibliography, images, or \input files.
% Numeric section labels and visible numbers are original UnsolvedMath IDs.
\documentclass[11pt,a4paper]{article}
\usepackage[margin=24mm,headheight=14pt]{geometry}
\usepackage{fontspec}
\setmainfont{Latin Modern Roman}
\setsansfont{Latin Modern Sans}
\setmonofont{Latin Modern Mono}
\usepackage{amsmath,amssymb,mathtools}
\usepackage{unicode-math}
\setmathfont{Latin Modern Math}
\usepackage{microtype}
\usepackage{enumitem,xurl,needspace}
\usepackage[dvipsnames]{xcolor}
\usepackage{hyperref}
\hypersetup{unicode=true,colorlinks=true,linkcolor=MidnightBlue,urlcolor=MidnightBlue,
 pdftitle={UnsolvedMath Problem 1920},pdfauthor={Source-annotated compilation},bookmarksnumbered=true}
\usepackage{bookmark,tocloft,titlesec,fancyhdr}
\setlength{\cftsecnumwidth}{4.2em}
\cftsetpnumwidth{2.5em}
\cftsetrmarg{4em}
\titleformat{\section}{\Large\bfseries\raggedright}{\thesection}{0.7em}{}
\pagestyle{fancy}\fancyhf{}
\fancyhead[L]{\small\sffamily Additional UnsolvedMath statements}
\fancyhead[R]{\small Original catalogue IDs}
\fancyfoot[C]{\thepage}
\setlength{\parindent}{0pt}
\setlength{\parskip}{5pt plus 1pt}
\setlength{\emergencystretch}{3em}
\setcounter{tocdepth}{1}\setcounter{secnumdepth}{1}
\setlist[itemize]{leftmargin=3em,itemsep=4pt,topsep=4pt,parsep=0pt}
\allowdisplaybreaks
\newcommand{\N}{\mathbb{N}}
\newcommand{\Z}{\mathbb{Z}}
\newcommand{\Q}{\mathbb{Q}}
\newcommand{\R}{\mathbb{R}}
\newcommand{\C}{\mathbb{C}}
\newcommand{\F}{\mathbb{F}}
\newcommand{\A}{\mathbb{A}}
\newcommand{\PP}{\mathbb{P}}
\newcommand{\E}{\mathbb{E}}
\newcommand{\eps}{\varepsilon}
\newcommand{\dd}{\,\mathrm d}
\newcommand{\sref}[2]{\hyperlink{source:#1:#2}{[S#2]}}
\newif\ifshowcatalogue
\showcataloguetrue
\newcommand{\cataloguescope}[1]{\ifshowcatalogue
 \paragraph{Statement recorded in the supplied source.}
 #1\fi}
\begin{document}
% UnsolvedMath ID 1920; source catalogue code EP-86

\setcounter{section}{1919}

\section{Four-Cycle-Free Subgraphs of the Hypercube}\label{1920}

{\small\sffamily UnsolvedMath ID 1920 · Graph Theory}

\subsection{Definitions and mathematical statement}

\paragraph{Shared notation.}Unless a section says otherwise, $\N=\{1,2,\ldots\}$,
$\Z$ is the ring of integers, $\Q,\R,\C$ are the rational, real, and complex
fields, $[N]=\{1,\ldots,\lfloor N\rfloor\}$, and $\log$ is the natural
logarithm. For real functions of a parameter tending to infinity,
$f\ll g$ and $f=O(g)$ mean $|f|\leq Cg$ eventually for some fixed $C$;
$f\gg g$ reverses this comparison; $f=o(g)$ means $f/g\to0$;
$f\sim g$ means $f/g\to1$; and $f\asymp g$ means both order bounds.
A subscript on $O$ or $\ll$ specifies allowed constant dependence.
The expression $n^{o(1)}$ in an upper bound means growth smaller than
$n^\epsilon$ for each fixed $\epsilon>0$. Local definitions govern any
exception, finite-field convention, or use of the same symbol for another
object. An asymptotic claim is not automatically an effective algorithm.



The vertices of $Q_n$ are $\{0,1\}^n$, adjacent when they differ in exactly one coordinate. Write $e_4(Q_n)$ for the maximum edges in a $C_4$-free subgraph. The asymptotic threshold question is whether $e_4(Q_n)/(n2^{n-1})\to1/2$; the source's $o(1)$ threshold is understood in this extremal sense.

A graph has a vertex set $V$ and edges that are unordered pairs of distinct vertices. A subgraph need not be induced unless that word is stated. A proper vertex colouring gives adjacent vertices different colours; $\chi(G)$ is the least number of colours.

$K_t$ is a complete graph and $C_t$ a simple cycle of length $t$. A complete multipartite graph has no edges within parts and all edges between different parts. An independent set has no internal edge; a clique has all its possible internal edges. \sref{1920}{1} \sref{1920}{2}

\paragraph{Statement recorded in the supplied source.}
Let $Q_n$ be the $n$-dimensional hypercube graph (so that $Q_n$ has $2^n$ vertices and $n2^{n-1}$ edges). Is it true that every subgraph of $Q_n$ with $ \geq \left(\frac{1}{2}+o(1)\right)n2^{n-1} $ many edges contains a $C_4$? \sref{1920}{1}

\subsection{Short English statement}

Is one half asymptotically the largest fraction of hypercube edges that can be retained without creating a four-cycle?

\subsection{Sources}

{\small

\begin{itemize}

\item[\hypertarget{source:1920:1}{S1}] ULAM.AI, \emph{UnsolvedMath}, supplied \texttt{problems.json}, record 1920 (EP-86), statement and background. The embedded literature-review date is 2026-08-17; that is the archive’s review date, not a fresh certification. Dataset landing page: \url{https://huggingface.co/datasets/ulamai/UnsolvedMath}.

\item[\hypertarget{source:1920:2}{S2}] T. F. Bloom, \emph{Erdős Problems}, Problem 86, \url{https://www.erdosproblems.com/86}. Reference imported from the supplied record; the live problem page was not independently checked for this compilation.

\item[\hypertarget{source:1920:3}{S3}] Rahil Baber, Turan densities of hypercubes, arXiv:1201.3587. \url{https://arxiv.org/abs/1201.3587}. Bibliographic information imported from the supplied record.

\item[\hypertarget{source:1920:4}{S4}] Jozsef Balogh, Ping Hu, Bernard Lidicky, and Hong Liu, Upper bounds on the size of 4- and 6-cycle-free subgraphs of the hypercube, European Journal of Combinatorics 35 (2014), 75-85. \url{https://arxiv.org/abs/1201.0209}. Bibliographic information imported from the supplied record.

\item[\hypertarget{source:1920:5}{S5}] Minamo Minamoto, New Lower Bounds for C4-Free Subgraphs of the Hypercubes Q7 and Q8, arXiv:2603.29127. \url{https://arxiv.org/abs/2603.29127}. Bibliographic information imported from the supplied record.

\item[\hypertarget{source:1920:6}{S6}] Brass, Peter and Harborth, Heiko and Nienborg, Hauke, On the maximum number of edges in a {$C_4$}-free subgraph of {$Q_n$}. J. Graph Theory (1995), 17--23. Bibliographic information imported from the supplied record.

\end{itemize}

}

\label{end:1920}
\end{document}
